\section{Numerical Methods}\label{sec:numerical}

To evaluate the model's quantitative performance we carry out two evaluations. First, as a proof of concept we analyze the impact of a set of reforms in Argentina from 2005 to 2010 that introduced near-universal pensions benefiting a large share of informal households (section \ref{sec:argentina}). Second, we calibrate the model to three developed economies for which almost all households are covered by contributive pensions and have different pension characteristics, aggregate spending, demographics, and income (and voting) inequality: the United States, France, and the United Kingdom (section \ref{sec:oecd}). Section \ref{sec:esc} then lets the electorate choose the design of the system as well as its size. This section describes how the equilibria behind those results are computed. The quantitative politico-economic literature typically computes a time-invariant policy function as the fixed point of a functional equation, by linear-quadratic approximation in \textcite{KrusellRR99a1} or by projection in \textcite{Song12a1}, or obtains it in closed form \parencite{mkvss08}. We instead compute a sequence of date-specific policy functions by backward induction from a terminal period. The object is the same, the Markov equilibrium in the limit of finite-horizon economies, which is the class in which \textcite{mkvss08} establish uniqueness; the terminal condition selects it where a stationary fixed point can admit several policy functions \parencite{Forni05a1}, and the non-stationary demographic path enters directly, with no stationarity required. What is specific to this paper is the structure that keeps the state space small, and how the solution exploits it.

\smalltitle{What is solved.}
Outside the case of proposition \ref{prop:LOG:PEE}, the politico-economic equilibrium of definition \ref{def:PEE} has no closed form. We identify the sequence of tax policy functions by backward induction, solving each date's policy function on a grid of states given the continuation policy of the next date, and then simulate forward, solving for the path of economic equilibria with taxes read off the policy functions at the simulated states. The recursion is initialized in one of two ways. In the finite-horizon economy that terminates after $T$ periods, the terminal policy is available analytically. In the infinite-horizon economy that converges to a steady state after $T$ periods, the steady-state policy function is a fixed point in which the policy and continuation policy functions are the same time-autonomous function of the state. With a sufficiently large $T$ the two agree out of steady state, and the following sections report the finite-horizon solution.\footnote{With $t$ measuring 30-year time steps, we use projected population growth rates until 2070--2080 and then default to five time steps with steady state rates.}

\smalltitle{The state space is small.}
In a Markov equilibrium the political decision at $t$ could in principle depend on the whole distribution of savings that retirees carry into the period. Proposition \ref{prop:LOG:EE}, part iii, removes it: relative savings $s_{t-1,i}/s_{t-1}$ are closed-form functions of the current tax alone, so the policymaker evaluates them at any candidate $\tau_t$ without tracking them as states, whatever the number of household types. Appendix \ref{app:EE:CRRA} shows that this survives under CRRA preferences. What remains is at most two scalars. In the model without informal households, log preferences leave no endogenous state at all. With informal savers, the ratio $\iota_{t-1} \equiv s_{t-1,0}/s_{t-1}$ does not reduce, because it depends on $\tau_{t-1}$ through the previous period's savings decision (appendix \ref{app:PEE:LOG_s0}); it is the single state under log preferences. Under CRRA preferences the marginal indirect utilities involve consumption levels rather than log-derivatives alone, and the aggregate savings level $s_{t-1}$ becomes a second state.

\smalltitle{Roots on a grid.}
Policies are identified from the first order condition of the political problem, $z_t \equiv \partial \mathcal{W}_t/\partial \tau_t = 0$, rather than by maximizing $\mathcal{W}_t$ directly, because the first order condition is where the closed-form savings shares do their work. A root problem, however, loses the safeguards of the maximization it came from: a root may be a minimum, there may be several, and the optimum may sit at a corner where there is no root. We therefore evaluate $z_t$ on a grid $\mathcal{T}$ of candidate taxes on $[l,u]$, with $u<1$ strictly since $z_t$ diverges as $\tau_t \rightarrow 1$, and let $\hat{z}_t$ denote its piecewise-linear interpolant. Its integral $\hat{\mathcal{V}}_t(\tau) = \int_l^{\tau}\hat{z}_t(x)\,dx$ reconstructs the political objective up to a constant from the same evaluations, and its maximizer on $[l,u]$ lies in the finite set
\begin{align}\label{eq:num:candidates}
    \mathcal{C}_t = \left\lbrace l,u \right\rbrace \cup \left\lbrace \tau \in (l,u) : \hat{z}_t(\tau) = 0 \text{ and } \hat{z}_t \text{ crosses from above} \right\rbrace.
\end{align}
The selected tax is the candidate with the largest $\hat{\mathcal{V}}_t$. Corners and interior roots thus compete under one criterion, upward crossings are excluded as minima, and the reconstructed objective only has to rank well-separated candidates, a much weaker requirement than the accuracy demanded of the roots themselves. Terms of $z_t$ that depend on the continuation policy have no closed form and are differentiated numerically along the same grid, in the coordinate $\ln(1-\tau_t)$ so that the singular part of every profile is affine and reproduced exactly. One derivative is never taken numerically: that of the old formal households' consumption must hold $s_{t-1,i}/s_{t-1}$ fixed, as the policymaker does, and is taken from its closed form.

\smalltitle{Log preferences.}
Log preferences make the recursion cheap and, in the model without informal households, exact. There no state survives and $z_t$ depends on $\tau_t$ alone, so the $T$ periods decouple into independent scalar problems, each solved on the grid and polished to machine precision by a bracketed root step. With informal savers, $z_t$ depends on $\tau_{t+1}$ through the continuation policy $\tau^{t+1}(\iota_t)$, and a policy function of $\iota_{t-1}$ has to be identified at each date. The state is cheap to carry, since it enters the first order condition only through the old informal households' term,
\begin{align}\label{eq:num:decomposition}
    z_t\left(\tau_t, \iota_{t-1}\right) = \bar{z}_t\left(\tau_t\right) + \omega \gamma_0 \mu_0 \dfrac{A_t}{\iota_{t-1} + A_t \tau_t},
\end{align}
where $A_t$ is a function of parameters alone, the informal households' discounted pension per unit of tax and of aggregate savings, and $\bar{z}_t$ collects every other term. Evaluating $z_t$ on the whole state grid costs one pass over $\mathcal{T}$ plus a rank-one correction, and the fixed point that determines $\iota_t$ as a function of $\tau_t$ does not involve $\iota_{t-1}$, so it is solved once per date rather than once per state.

\smalltitle{CRRA preferences.}
Under CRRA preferences the policy function is defined over the savings level as well, and each date requires a fixed point: a candidate $s_t$ is the equilibrium at $(\tau_t, s_{t-1})$ when it reproduces itself through the savings decision, with the continuation policies evaluated at $s_t$. Two features of the model keep this affordable. First, the predetermined state enters the fixed point only through the factor $(s_{t-1}/\nu_t)^{\sigma_{s,t}}$ of appendix \ref{app:EE}, so the costly part is evaluated once over candidate taxes and candidate savings, and the residual for every $s_{t-1}$ follows by multiplying it by a scalar. Second, with informal savers the two states unnest exactly: the fixed point in $\iota_t$ involves neither predetermined state, so it is solved first as a function of $(\tau_t, s_t)$, and the fixed point in $s_t$ remains one-dimensional. Candidate taxes for which no equilibrium state lies inside the grid, or at which a consumption level is non-positive, are marked infeasible rather than clipped, and \eqref{eq:num:candidates} is applied to the feasible sub-grid. The selected policy functions are lightly smoothed before they are interpolated as continuation policies, with the smoother's knots fixed in advance, so that the map from parameters to solution is continuous; the calibration below relies on this. At $\rho = 1$ the CRRA recursion reproduces the log solution, and since the two solvers share almost no code, their agreement as $\rho \rightarrow 1$ is the main check on both.

\smalltitle{Path and calibration.}
The recursion delivers policy functions. To report an equilibrium we take the state entering the first period from the steady state under the first period's own policy, a scalar fixed point in $\tau_1$, walk the policy functions forward, and re-solve the economic equilibrium exactly at the resulting tax path, so that reported allocations carry the precision of the equilibrium conditions rather than of the grids.

Calibration is a nested fixed point. The inner problem is the full politico-economic equilibrium path at a trial parameter vector. The outer problem searches over the parameters that cannot be set without it until the targets of sections \ref{sec:argentina} and \ref{sec:oecd} hold at the calibration year: $(\beta, \omega)$ for the rich economies, targeting the interest rate and the tax rate, and $(\beta, \omega, \eta_0, X_0)$ for Argentina, targeting the capital-output ratio, the tax rate, and the informal households' relative income and hours. Every other parameter is set before the loop. Since $\beta$ moves the savings state entering the calibration year through every preceding period, the inner solve cannot be shortened to that year. Results across the intertemporal elasticity of substitution come from a march anchored at $\rho = 1$, the one value at which the cheap log solver applies, each further $\rho$ started from an extrapolation through the two nearest solved points. Along that march the calibrated $\eta_0$ and $X_0$ barely move while $\beta$ and $\omega$ move by tens of percent, so it is the latter pair that the data targets genuinely identify.

Section \ref{sec:esc} lets the electorate choose the design $\theta$ as well as the tax. That choice is layered on the same machinery with one difference: because whether the chosen design is interior is the question asked there, the political objective is evaluated directly on a grid of designs and maximized, so that a corner is observed rather than inferred from a root solver's failure. The full technical documentation of the algorithms, their calibration protocols, and the checks behind every structural claim made here is published alongside the paper, together with the code.\footnote{\textit{Social Security Design and Its Political Support: Technical Documentation}, and the repository at \url{https://github.com/ChampionApe/SSDclaude}.}
